\documentclass[11pt]{article}
\usepackage{coursenotes}
\begin{document}
\notesheader{3}{Conditional Effects}{Who responds, how to estimate it, and why outcome accuracy is the wrong test}

\begin{decision}
An offer works on average. Whom should it go to? The answer depends on how the effect varies with customer
characteristics, the conditional average treatment effect, held against a break-even that may itself vary
by customer. Models that predict outcomes well can still target badly.
\end{decision}

\begin{goals}
  \item define the CATE, identify it in an experiment, and state where it is not identified;
  \item turn a CATE into a targeting rule under per-unit and per-redemption costs;
  \item read an interaction model and test whether two segments really differ;
  \item explain how S-, T-, X-, DR- and R-learners estimate the CATE, and what each one risks;
  \item show why outcome prediction error cannot choose a targeting model, and name losses that can.
\end{goals}

%=====================================================================================
\sectionone

\subsection{The estimand}

With treatment $D$, outcome $Y$ and pre-treatment covariates $X$, write $\mu_d(x) = \E[Y \mid D=d, X=x]$ and
\[
  \tau(x) = \E[Y(1) - Y(0) \mid X = x],
\]
the \emph{conditional average treatment effect} (CATE). It is a function; the ATE is its average. Once
identified, every outcome regression splits into a \emph{baseline} and an \emph{effect}:
\begin{equation}\label{eq:decomp}
  \mu_d(x) = \mu_0(x) + d\,\tau(x).
\end{equation}
A decision about a target population $T$ (a list, a region) needs $\tau_T = \sum_x \Pr_T(X=x)\,\tau(x)$:
standardisation with $T$'s weights.

A variable is \emph{prognostic} if $\mu_0(x)$ varies with it and an \emph{effect modifier} if $\tau(x)$ does.
The two are independent properties: a variable can be either, both or neither. In most business data
$\Var(\mu_0(X))$ is far larger than $\Var(\tau(X))$, which is why learners tuned to outcomes favour
prognostic variables.

\subsection{Identification}

\begin{assum}{Conditional independence}{ci}
$\{Y(0), Y(1)\} \indep D \mid X$.
\end{assum}
\begin{assum}{Overlap on the decision population}{ov}
$0 < \Pr(D=1 \mid X=x) < 1$ for every $x$ at which a decision will be made.
\end{assum}
\begin{prop}{Identification of the CATE}{id}
Under consistency and Assumptions~\ref{asm:ci}--\ref{asm:ov}, $\tau(x) = \mu_1(x) - \mu_0(x)$.
\end{prop}
\begin{proof}
$\mu_d(x) = \E[Y(d) \mid D=d, X=x]$ by consistency (well defined by overlap) $= \E[Y(d) \mid X=x]$ by
conditional independence. Subtract.
\end{proof}

An experiment satisfies conditional independence by design, and overlap for every $x$ \emph{in the
experiment}. A decision population with values of $x$ the experiment never saw (new customers, a new
region) has no identified CATE there; any estimate is the model's extrapolation. Keep the two failures
apart: if treatment is a deterministic function of $X$, conditional independence holds trivially but overlap
fails everywhere (a support problem); if treatment depends on something outside $X$, overlap may hold while
conditional independence fails (a confounding problem).

\subsection{From the CATE to a targeting rule}

Let a treated unit generate margin $m$ per purchase, cost $c$ per purchase made while treated (a redeemed
coupon) and $k$ per unit treated regardless (a phone call).

\begin{prop}{Individual targeting rule}{rule}
The expected net value of treating a unit with covariates $x$ is $V(x) = m\,\tau(x) - c\,\mu_1(x) - k$.
Without a budget, the value-maximising rule treats $x$ if and only if $V(x) > 0$, and its value per unit is
\[
  \E\big[\max\{V(X), 0\}\big] \;\ge\; \max\{\E[V(X)],\, 0\},
\]
the value of the best blanket rule. With a budget of $B$ treatments, treat the $B$ units with the largest
positive $V(x)$.
\end{prop}
\begin{proof}
Treating minus not treating is $\E[(m-c)Y(1) - k \mid x] - \E[mY(0) \mid x]$, which is $V(x)$ by
\eqref{eq:decomp}. Total value is a sum over units, so each is treated iff its term is positive. The inequality
holds because $\max\{V,0\} \ge V$ and $\max\{V,0\} \ge 0$ pointwise; take expectations.
\end{proof}

The gap in the inequality is the \emph{value of targeting}. It is large when effects straddle the break-even:
an average below break-even can hide a profitable segment. With per-unit costs ($c=0$) the rule ranks by
$\tau(x)$ against a common threshold $k/m$. With per-redemption costs ($k=0$) the threshold $c\,\mu_1(x)/m$
rises with purchase rates, so likely buyers must clear a higher bar. Either way, ranking by ``most likely to
buy'' is not the rule and can be its opposite: in the four response types of Meeting~1, a response model ranks
sure things first, while the rule wants persuadables.

\subsection{The interaction model and inference on heterogeneity}

With low-dimensional $X$, the model $Y = \alpha + \tau D + \beta^\top X + \gamma^\top (D \cdot X) + \varepsilon$ gives
$\mu_0(x) = \alpha + \beta^\top x$ and $\tau(x) = \tau + \gamma^\top x$. With one binary $X$ it reproduces the four
cell means exactly. Its fitted values give every row both $\hat\mu_0(X_i)$ and $\hat\mu_1(X_i)$, whichever arm it
was in. Two cautions: $\tau$ is the effect at the reference group $x=0$ and recoding the reference changes its
value and the sign of $\gamma$, but not the fitted effects; and the average effect for $T$ is
$\tau + \gamma^\top\E_T[X]$, not $\tau$.

\begin{prop}{Standard error of a difference between segments}{gap}
For segment effects estimated on disjoint randomised samples,
$\SE(\hatt_a - \hatt_b) = \sqrt{\SE(\hatt_a)^2 + \SE(\hatt_b)^2}$.
\end{prop}
\begin{proof}
Disjoint samples make the two estimates independent, and variances of independent terms add.
\end{proof}

Heterogeneity is shown by an interval for the \emph{difference} that excludes zero, not by one segment being
``significant'' and another not. With $K$ pre-specified segments each tested at level $\alpha$, the chance of a
false positive somewhere approaches $1 - (1-\alpha)^K$; report all segments and, if the decision rests on
several, use Bonferroni ($\alpha/K$) or Holm. Segments found by searching need fresh data (Meeting~4). The
segment-level effects $\text{GATE}_k = \E[\tau(X) \mid G=k]$ for a grouping $G$ are estimated by within-group
differences in means.

\subsection{Meta-learners}

Meta-learners turn any prediction model into a CATE estimator.

\paragraph{S-learner.} One model $\hat f(x, d)$ with $D$ as a feature; $\hatt(x) = \hat f(x,1) - \hat f(x,0)$.

\paragraph{T-learner.} $\hat\mu_1$ fitted on treated rows, $\hat\mu_0$ on control rows;
$\hatt(x) = \hat\mu_1(x) - \hat\mu_0(x)$.

Both minimise a loss in $Y$ and get $\hatt$ as a by-product. The next two results say what that costs.

\begin{prop}{Regularisation shrinks heterogeneity in the S-learner}{shrink}
In a balanced design with centred, orthogonal regressors $\tilde X$, $\tilde D = D - \tfrac12$ and
$\tilde D\tilde X$, penalising only the interaction by $\lambda\gamma^2$ gives
$\hat\gamma_\lambda = \hat\gamma_{\text{OLS}} \cdot s/(s + \lambda)$ with $s = \sum_i (\tilde D_i \tilde X_i)^2$; the other
coefficients are unchanged.
\end{prop}
\begin{proof}
Orthogonality separates the objective by coefficient; the $\gamma$ part, $\sum_i (Y_i - \gamma W_i)^2 + \lambda\gamma^2$
with $W_i = \tilde D_i\tilde X_i$, is minimised at $\sum_i W_iY_i/(s+\lambda)$.
\end{proof}
Since $\tilde D = \pm\tfrac12$, $s = \tfrac14\sum_i \tilde X_i^2$: the interaction column has a quarter of $X$'s
variation, so a penalty that barely moves $\beta$ can erase $\gamma$. A lasso sets a small $\gamma$ exactly to zero,
and the model then reports the same effect for everyone. Trees do the same implicitly: a split on $D$ competes
with prognostic splits and loses when the effect is small.

\begin{prop}{Variances add in the T-learner}{tlearner}
If $\hat\mu_1$ and $\hat\mu_0$ are fitted on disjoint samples,
$\Var(\hatt(x)) = \Var(\hat\mu_1(x)) + \Var(\hat\mu_0(x))$ and
$\text{bias}(\hatt(x)) = \text{bias}(\hat\mu_1(x)) - \text{bias}(\hat\mu_0(x))$.
\end{prop}
\begin{proof}
The estimates are independent, so their variances add; expectation is linear, so their biases subtract.
\end{proof}
Variances always add; biases can cancel or reinforce. By \eqref{eq:decomp} each arm must learn the whole
baseline, usually most of the signal. Fitted on different halves with different regularisation, the two arms
smooth the baseline differently, and the difference lands in $\hatt$. With truth $\mu_1 = 0.30$, $\mu_0 = 0.10$,
fits $0.34$ and $0.14$ miss each outcome by $0.04$ but get $\tau = 0.20$ exactly; fits $0.32$ and $0.08$ miss by
only $0.02$ but give $0.24$. For the effect, what matters is whether the two errors move together, not how
large each one is.

\paragraph{X-learner.} Fit $\hat\mu_0, \hat\mu_1$; impute $\tilde\tau_i = Y_i - \hat\mu_0(X_i)$ for treated units and
$\tilde\tau_i = \hat\mu_1(X_i) - Y_i$ for controls; fit $\hatt_1$ and $\hatt_0$ to the two sets of imputations; combine
$\hatt(x) = e(x)\,\hatt_0(x) + (1 - e(x))\,\hatt_1(x)$. Each imputation uses the \emph{other} arm's model, so the
larger arm's model does most of the work; it helps when arms are unbalanced.

\paragraph{Learners that target the effect.} The next three learners regress something whose conditional
mean \emph{is} $\tau(x)$, so their loss is about the effect rather than the outcome.

\begin{prop}{The transformed outcome}{tro}
With known propensity $e(x)$, let $Y^\ast = Y\,\dfrac{D - e(X)}{e(X)(1-e(X))}$. Under Assumptions~\ref{asm:ci}
and \ref{asm:ov}, $\E[Y^\ast \mid X] = \tau(X)$.
\end{prop}
\begin{proof}
$Y^\ast = DY/e(X) - (1-D)Y/(1-e(X))$. Given $X$, $\E[DY \mid X] = e(X)\mu_1(X)$ and
$\E[(1-D)Y \mid X] = (1-e(X))\mu_0(X)$; divide and use Result~\ref{prop:id}.
\end{proof}
Any regression of $Y^\ast$ on $X$ estimates the CATE, but each $Y^\ast_i$ is one outcome scaled up, so it is very
noisy.

\begin{prop}{The doubly robust pseudo-outcome}{dr}
Let
\[
  \phi = \hat\mu_1(X) - \hat\mu_0(X) + \frac{D\,(Y - \hat\mu_1(X))}{\hat e(X)} - \frac{(1-D)(Y - \hat\mu_0(X))}{1 - \hat e(X)},
\]
with $\hat\mu_d$ and $\hat e$ fitted on other data. Then $\E[\phi \mid X] = \tau(X)$ if \emph{either} $\hat\mu_0 = \mu_0$
and $\hat\mu_1 = \mu_1$, \emph{or} $\hat e = e$.
\end{prop}
\begin{proof}
Given $X$, $\E[D(Y - \hat\mu_1)/\hat e \mid X] = e(\mu_1 - \hat\mu_1)/\hat e$, and similarly for the control term.
So $\E[\phi \mid X] - \tau = (\mu_1 - \hat\mu_1)(e/\hat e - 1) - (\mu_0 - \hat\mu_0)\big((1-e)/(1-\hat e) - 1\big)$,
which is zero if the outcome models are right or the propensity is right.
\end{proof}
The \emph{DR-learner} regresses $\phi$ on $X$ (with cross-fitting). The outcome models remove most of the noise in
$Y^\ast$; the weighted residuals correct their bias. In an experiment $e$ is known, so $\phi$ is unbiased even with
poor outcome models.

\begin{prop}{The R-learner's target}{rlearner}
Let $\ell(x) = \E[Y \mid X=x]$. Under Assumption~\ref{asm:ci}, $Y - \ell(X) = (D - e(X))\,\tau(X) + \varepsilon$ with
$\E[\varepsilon \mid D, X] = 0$, so $\tau$ minimises $\E\big[\{(Y - \ell(X)) - (D - e(X))\,t(X)\}^2\big]$ over functions $t$.
\end{prop}
\begin{proof}
Write $Y = \mu_0(X) + D\tau(X) + \varepsilon$ with $\E[\varepsilon \mid D, X] = 0$. Then
$\ell(X) = \mu_0(X) + e(X)\tau(X)$; subtract. The squared loss is minimised by the conditional mean, and
the residual is orthogonal to any function of $(D, X)$.
\end{proof}
The \emph{R-learner} minimises the sample version, the \emph{R-loss}, with $\hat\ell$ and $\hat e$ fitted on other
folds and a penalty on $t$. It regresses the outcome residual on the treatment residual, locally in $x$: the same
idea as causal forests (Meeting~4) and double machine learning (Meeting~8).

\paragraph{A shared representation.} Neural networks offer a middle ground between sharing everything
(S-learner) and nothing (T-learner): TARNet passes $x$ through shared layers to a representation $\Phi(x)$ and
then through two heads, $\hat\mu_d(x) = h_d(\Phi(x))$, training each unit only on its own arm's head. Variants
that force $\Phi(X)$ to look alike in both arms can hide an overlap failure rather than fix it. Handout H6 covers
these architectures.

\subsection{Why outcome accuracy cannot choose a targeting model}

\begin{prop}{The cost of ignoring all heterogeneity}{blind}
Under a 50/50 design, the best predictor of $Y$ that imposes a constant effect, $g(x) + \bar\tau d$, has mean
squared error exceeding that of the correct model $\mu_D(X)$ by exactly $\tfrac14\Var(\tau(X))$.
\end{prop}
\begin{proof}
For fixed $\bar\tau$ the best $g$ at each $x$ is $\mu_0(x) + \tfrac12(\tau(x) - \bar\tau)$, leaving errors
$\pm\tfrac12(\tau(x) - \bar\tau)$ in the two arms. Their average square is minimised at $\bar\tau = \E[\tau(X)]$,
giving $\tfrac14\Var(\tau(X))$; the irreducible noise is common to both models.
\end{proof}

A quarter of the effect's variance is tiny next to an outcome's own variance (about $p(1-p)$ for a purchase).
A model can be nearly optimal by outcome loss while discarding all the heterogeneity a targeting decision
needs. Since $\tau_i$ is never observed, the error we want, $\E[(\hatt(X) - \tau(X))^2]$, is computable only in
simulations. On a held-out randomised fold, compare candidates by (i) the value of the policy each implies
(Meeting~5); (ii) the R-loss; or (iii) the mean squared gap between $\hatt(X)$ and the DR pseudo-outcome $\phi$.
When nuisance models are accurate, (ii) and (iii) rank candidates as their true CATE error would. Assign
\emph{units}, not rows, to the samples used for fitting, choosing and assessing, so that no customer appears in two.

\begin{pitfall}[targeting errors]
Ranking by predicted purchase; declaring heterogeneity because one segment is significant; choosing a CATE
model by outcome RMSE; reporting the effect at the reference group as ``the'' effect.
\end{pitfall}

\begin{center}
\begin{tikzpicture}
\begin{axis}[width=0.72\linewidth, height=5cm, domain=0:10, samples=60, xlabel={tenure (months)},
  ylabel={renewal probability}, ymin=0.3, ymax=1, axis lines=left,
  legend style={draw=none, font=\footnotesize, at={(0.02,0.98)}, anchor=north west}]
\addplot[thick, extgrey] {0.45 + 0.045*x};
\addplot[thick, sbgreen] {0.45 + 0.045*x + 0.22 - 0.02*x};
\legend{$\mu_0(x)$: without offer, $\mu_1(x)$: with offer}
\end{axis}
\end{tikzpicture}

\small Figure 1. Tenure is strongly prognostic (both curves rise) and also an effect modifier (the gap shrinks).
A response model sees the rise; targeting needs the gap.
\end{center}

%=====================================================================================
\sectiontwo

\paragraph{Setting.} FitLife tests a retention offer (a free month, costing ¥80 per member offered) on members
near the end of their contracts. The outcome is renewal; a renewal is worth ¥600 of margin, so the break-even is
$\tau^\ast = 80/600 = 0.133$. Members are split by tenure (short: under a year; long) and usage (low; high) into
four segments of 1{,}000, each randomised 50/50.

\paragraph{Step 1: segment effects.}
\begin{center}
\begin{tabular}{lcccccc}
\toprule
Segment & Offered & Not offered & $\hatt_k$ & SE & 95\% CI & $z$ vs $\tau^\ast$ \\
\midrule
Short, low  & 0.62 & 0.50 & 0.12 & 0.031 & $[0.059, 0.181]$ & $-0.43$ \\
Short, high & 0.70 & 0.48 & 0.22 & 0.030 & $[0.161, 0.279]$ & $2.86$ \\
Long, low   & 0.84 & 0.80 & 0.04 & 0.024 & $[-0.008, 0.088]$ & $-3.85$ \\
Long, high  & 0.93 & 0.90 & 0.03 & 0.018 & $[-0.005, 0.065]$ & $-5.87$ \\
\bottomrule
\end{tabular}
\end{center}
Each SE is $\sqrt{p_1(1-p_1)/500 + p_0(1-p_0)/500}$. The ATE is $0.1025$: offering it to everyone nets
$600 \times 0.1025 - 80 = -$¥18.50 per member. Net values by segment are $-8$, $+52$, $-56$ and $-62$, so the
targeted rule (short, high usage only) earns $52/4 = $ ¥13 per member in the base against $0$ for the best
blanket rule (offer to nobody): the value of targeting in Result~\ref{prop:rule}. With four segments, Bonferroni
at $0.0125$ still accepts the short, high-usage segment ($p = 0.002$).

\paragraph{Step 2: prognostic versus modifier.} Tenure is strongly prognostic (control renewal $0.49$ for short,
$0.85$ for long) and a modifier. Usage barely predicts renewal among short-tenure members ($0.50$ vs $0.48$) but
modifies the effect ($0.12$ vs $0.22$). A churn model ranking by low predicted renewal treats the two short-tenure
segments alike and so includes one that loses ¥8 per offer.

\paragraph{Step 3: an S-learner without interactions.} $Y = \alpha + \beta D + \gamma_1\,\text{long} + \gamma_2\,\text{high}$
gives $\hat\beta = 0.1025$ and hence the same $\hatt(x)$ for every member: it recommends offering to nobody.

\paragraph{Step 4: a T-learner with additive arms.} Fitting $Y = a_d + b_d\,\text{long} + c_d\,\text{high}$ in each arm
gives $(0.6175, 0.225, 0.085)$ for offered and $(0.47, 0.36, 0.04)$ for not offered, and implied effects
$0.148$, $0.193$, $0.013$, $0.058$ against true $0.12$, $0.22$, $0.04$, $0.03$. It places short, low-usage members above
break-even ($0.148 > 0.133$) because an additive model in each arm cannot represent the tenure $\times$ usage
interaction in the effect. (A T-learner and an S-learner that interacts $D$ with every feature are the same
model.)

\paragraph{Step 5: the transformed outcome.} With $e = 0.5$, $Y^\ast = 2Y$ for offered and $-2Y$ for others. Every
value is $-2$, $0$ or $2$, yet its segment mean is the segment effect: for short, high usage,
$0.5 \times 2 \times 0.70 - 0.5 \times 2 \times 0.48 = 0.22$ (Result~\ref{prop:tro}).

\begin{exercise}[computation]
Is the difference between the two short-tenure segments ($0.22$ vs $0.12$) statistically clear? Compute its SE
and $z$.
\end{exercise}
\answer{$\SE = \sqrt{0.0303^2 + 0.0312^2} = 0.0435$; $z = 0.10/0.0435 = 2.30$, two-sided $p \approx 0.02$. The two
segments differ; it is this test, not the two separate significance levels, that shows heterogeneity.}

\begin{exercise}[judgement]
FitLife's data team reports that its gradient-boosted churn model has the lowest validation log-loss of all
candidates and proposes to send the offer to the 25\% of members it rates most likely to leave. Give two
reasons this could lose money, and the comparison you would ask for instead.
\end{exercise}
\answer{(i) Log-loss rewards predicting renewal, dominated by prognostic variables (Result~\ref{prop:blind}); it
says almost nothing about who \emph{responds}. (ii) The members most likely to leave include lost causes and the
short, low-usage segment, whose effect is below break-even. Ask for each candidate's policy value on a held-out
randomised fold (net of the ¥80 cost), or its R-loss or DR-pseudo-outcome loss on that fold.}

%=====================================================================================
\sectionthree

\begin{extension}[Theory of meta-learners]
K\"unzel et al.\ (2019) show when the X-learner beats the T-learner (notably with unbalanced arms). Nie and Wager
(2021) prove that the R-learner attains the rate of an oracle that knew $\ell$ and $e$, provided those are learned
fast enough: the property called Neyman orthogonality. Kennedy (2023) gives the analogous result for the
DR-learner.
\end{extension}

\begin{extension}[Lasso for heterogeneity]
With many candidate modifiers, the modified-outcome or R-loss lasso selects a sparse $\tau(x)$. Selection makes
naive standard errors invalid; use sample splitting or the best-linear-projection test of Meeting~4.
\end{extension}

\subsection*{Annotated reading}
\begin{readinglist}
\reading{K\"unzel, Sekhon, Bickel and Yu (2019), ``Metalearners for estimating heterogeneous treatment
effects'', \emph{PNAS}}{The S-, T- and X-learners compared, with simulations}{Sections 1--4}{2}
\reading{Nie and Wager (2021), ``Quasi-oracle estimation of heterogeneous treatment effects'',
\emph{Biometrika}}{The R-learner and why residualising first protects the effect}{Sections 1--3}{3}
\reading{Kennedy (2023), ``Towards optimal doubly robust estimation of heterogeneous causal effects'',
\emph{Electronic Journal of Statistics}}{The DR-learner and its error bound}{Sections 1--3}{3}
\reading{Gelman and Stern (2006), ``The difference between significant and not significant is not itself
statistically significant'', \emph{American Statistician}}{The segment-comparison error in four pages}{All}{1}
\reading{Handout H6, \emph{Neural and foundation-model estimators}}{TARNet, DragonNet, CEVAE and in-context
estimators, with what each assumes}{All}{2}
\end{readinglist}

\end{document}
